\documentclass[../main.tex]{subfiles}

\begin{document}
\begin{problema}
    Sea \(f \colon \mathbb{C} \to \mathbb{C}\). Pruebe que

    \begin{enumerate}
        \item Si \(f\) es continua en 0 y \(f(z) = f(2z)\) para toda \(z\in
            \mathbb{C}\),
            entonces \(f\) es constante.
        \item Si \(f\) es analítica en \(\mathbb{C}\), \(f^{\prime}\) es
            continua en 0, y
            \(f(2z) = 2f(z)\) para toda \(z\in \mathbb{C}\), entonces existe \(c\in
            \mathbb{C}\) tal que \(f(z) = cz\) para toda \(z\in \mathbb{C}\).
    \end{enumerate}
\end{problema}

\startsolution

\begin{enumerate}
    \item
        Notemos que

        \begin{equation*}
            f(z) = f\bigl(\dfrac{1}{2^{n}}z\bigr).
        \end{equation*}

        Tomando \(n \to \infty\) obtenemos que

        \begin{equation*}
            f(z) = f(0),\quad \forall z\in \mathbb{C}.
        \end{equation*}

        Por lo que \(f\) es constante.

    \item
        Derivando la igualdad que nos dan

        \begin{align*}
            \bigl(f(2z)\bigr)^{\prime} &= \bigl(2f(z)\bigr)^{\prime},\\
            2f^{\prime}(2z) &= 2f^{\prime}(z),\\
            f^{\prime}(2z) &= f^{\prime}(z).
        \end{align*}

        Y como \(f^{\prime}\) es continua en 0, tenemos que la igualdad
        anterior se satisface si

        \begin{align*}
            f^{\prime}(z) &= c,\\
            \implies f(z) &= cz + d.
        \end{align*}

        Por lo que,

        \begin{align*}
            f(2z) &= 2f(z),\\
            c(2z) + d &= 2(cz + d),\\
            2cz - 2cz - 2d + d &= 0,\\
            \therefore d &= 0.
        \end{align*}

        Y por lo tanto, \(f(z) = cz,\medspace \forall z\in \mathbb{C}\).

\end{enumerate}
\end{document}
